\LectureAnswerSection

\begin{drillsolution}{k6p15}\textbf{在实际斜绳上求杆端速度的投影}\par\nobreak
投影法：设图\ref{k6p15q}的 $A=(x,y)$，则杆端速度为 $(\omega y,-\omega x)$，沿 $B\to A$的单位方向为 $(x,y-h)/\ell$。点积逐项相乘，斜段增长率为
\[
 \bm v_A\cdot\bm e_{BA}
 =\frac{\omega xy-\omega x(y-h)}\ell
 =\frac{\omega hx}{\ell}=\omega h\sin\alpha.
\]
最后一步使用 $x/\ell=\sin\alpha$。悬挂段相应缩短，悬物向上。
\fig{k6p15s}{同一斜绳上，速度投影对应长度的微小变化}{\begin{tikzpicture}
\coordinate(O)at(0,0);\coordinate(B)at(0,2.6);\coordinate(A)at(2.4,.8);
\draw[world](O)--(A);\draw[world](B)--(A);\draw[guide](O)--(B)node[midway,left]{$h$};
\foreach \n in {O,A,B}{\fill(\n)circle(1.4pt);}\node[left]at(O){$O$};\node[above]at(B){$B$};\node[right]at(A){$A$};
\coordinate(V)at(2.72,-.16);\coordinate(H)at(3.0656,.3008);
\draw[flow,blue](A)--(V)node[below]{$\bm v_A\Delta t$};\draw[flow,red](A)--(H)node[right]{$\Delta\ell$};\draw[guide](V)--(H);
\draw[guide](A)--(3.4,.05);\node at(1.3,1.9){$\ell$};
\end{tikzpicture}}



\par\medskip\Needspace{4\baselineskip}\noindent\textbf{对同一几何长度微分并核对正负号}\par\nobreak
绳长法：$\ell^2=x^2+(y-h)^2$，$h$恒定，链式求导得
\[
 2\ell\dot\ell=2x\dot x+2(y-h)\dot y.
\]
除以正的 $2\ell$，再代入 $\dot x=\omega y,\dot y=-\omega x$，
\[
 \dot\ell=\frac{x(\omega y)+(y-h)(-\omega x)}\ell
 =\frac{\omega xy-\omega xy+\omega hx}{\ell}
 =\frac{\omega hx}{\ell}.
\]
因此 $\boxed{v_{\text{上}}=\omega h\sin\alpha}$。图\ref{k6p15s}把杆端微小位移沿 $BA$投影，红色投影就是斜绳增长量；用坐标展开这一投影便得到长度微分式，二者可互查正负号。若转向改为逆时针，$\dot x,\dot y$同时反号，悬物改为向下同速，不能只更改图上箭头而保留“上升”结论。

两种方法的共同基础是距离导数 $\dot\ell=\bm e_{BA}\cdot(\bm v_A-\bm v_B)$，本题 $B$固定。用坐标求长度微分只是将同一点积展开，速度投影则直接读出它的几何意义。悬挂段若用向下坐标，便有 $\dot y_M=-\dot\ell$；改用向上坐标后符号再反转，才得到正的上升速率。这一步统一了方程正号与图中箭头。\end{drillsolution}

\begin{drillsolution}{k6p16}\textbf{固定横坐标后优化发射角参数}\par\nobreak
对图\ref{k6p16q}的 $x\ne0$，令 $p=\tan\alpha$，轨迹为 $y=xp-gx^2(1+p^2)/(2v_0^2)$。在固定 $x$处对 $p$求导时，$x,v_0,g$均为常量：
\begin{align*}
 \frac{\partial y}{\partial p}
 &=x-\frac{gx^2}{2v_0^2}(2p)
 =x-\frac{gx^2}{v_0^2}p,\\
 \frac{\partial^2y}{\partial p^2}&=-\frac{gx^2}{v_0^2}<0.
\end{align*}
令一阶导数为零，得到唯一最大值的参数 $p_*=v_0^2/(gx)$。代回高度时分别计算一次项、平方项和常数项：
\begin{align*}
 y(p_*)&=x\frac{v_0^2}{gx}
 -\frac{gx^2}{2v_0^2}
 -\frac{gx^2}{2v_0^2}\frac{v_0^4}{g^2x^2}\\
 &=\frac{v_0^2}{g}-\frac{gx^2}{2v_0^2}-\frac{v_0^2}{2g}.
\end{align*}
合并同类项，得到
\[\boxed{y_{\rm env}=\frac{v_0^2}{2g}-\frac{gx^2}{2v_0^2}}.\]
\fig{k6p16s}{给定初速的抛射包络：解答过程}{\begin{tikzpicture}\draw[axis](0,0)--(5,0)node[right]{$p=\tan\alpha$};\draw[axis](0,0)--(0,3)node[above]{$y(x,p)$};\draw[world]plot[domain=.1:4.2,samples=50]({\x},{2.5-.5*(\x-2.2)^2});\draw[guide](2.2,0)--(2.2,2.5);\node[below]at(2.2,0){$v_0^2/(gx)$};\node[above]at(2.2,2.5){包络高度};\end{tikzpicture}}



\par\medskip\Needspace{4\baselineskip}\noindent\textbf{区分角度参数曲线与空间包络}\par\nobreak
图\ref{k6p16s}的横轴是发射角参数 $p$，它画的是固定 $x$ 处高度随角度的变化；顶点的纵坐标才给这一 $x$ 处的包络高度。把各 $x$ 对应的顶点高度连起来，得到题设图\ref{k6p16q}中的蓝色空间包络，每条发射轨迹与它相切。包络下方的点通常对应判别式为正的两种发射角，包络上为一个临界角，上方不可达。$x=0$需单独使用竖直抛射，最高可达 $v_0^2/(2g)$，与公式连续。若另有地面终止条件，只取落地前实际经过的部分；不能把地面以下的数学延长线当成飞行路径。



\par\medskip\Needspace{4\baselineskip}\noindent\textbf{用配方验证上界并恢复实际时间方向}\par\nobreak
还可在固定 $x\ne0$处记 $K=gx^2/(2v_0^2)>0$，将高度配方：
\begin{align*}
 y&=-K\left[p^2-\frac xKp+\frac{x^2}{4K^2}\right]
 +\frac{x^2}{4K}-K\\
 &=y_{\rm env}-K\left(p-\frac{v_0^2}{gx}\right)^2.
\end{align*}
平方项非负，因此每一轨迹都在包络下方；等号给相切位置对应的发射角。对地面以上、$x>0$的点，两实根都应选 $\cos\alpha>0$的代表角，保证正时间；若仅使用正切而忘记角度所在象限，可能把指向目标反方向的角也重复计入。\end{drillsolution}

\begin{drillsolution}{k6p17}\textbf{用跑游时间比例建立平均速度}\par\nobreak
取图\ref{k6p17q}的岸线为 $x$轴，向水面为 $y$正向。若在总时间 $T$后追上，跑步占比为 $\lambda\in[0,1)$，游泳方向为单位矢量 $\bm e$，两者总位移相等给
\[UT(\cos15^\circ,\sin15^\circ)=4\lambda T\bm e_x+2(1-\lambda)T\bm e.
\]


\par\medskip\Needspace{4\baselineskip}\noindent\textbf{沿支持线法向证明所有方案的速度上界}\par\nobreak
由于两人同时同地出发，可以约去 $T>0$。要对所有方案给出共同上界，取图\ref{k6p17s}中的单位方向 $\bm n=(1/2,\sqrt3/2)$。跑步速度在该方向投影为2，任何游泳速度的投影不超过2，所以
\[U\cos45^\circ
=4\lambda\bm e_x\cdot\bm n+2(1-\lambda)\bm e\cdot\bm n
\le2\lambda+2(1-\lambda)=2.
\]
\fig{k6p17s}{先跑后游追赶离岸小船：解答过程}{\begin{tikzpicture}\draw[axis](0,0)--(5,0)node[right]{$V_x$};\draw[axis](0,0)--(0,3.2)node[above]{$V_y$};\draw[world](0,0)circle(2);\draw[world](4,0)--(1,1.732);\draw[ray](0,0)--(3.7,.991);\fill(2.732,.732)circle(2pt)node[above right]{临界船速};\fill(4,0)circle(2pt)node[below]{$4$};\fill(1,1.732)circle(2pt)node[above]{$(1,\sqrt3)$};\end{tikzpicture}}

因此 $U\le2\sqrt2$。



\par\medskip\Needspace{4\baselineskip}\noindent\textbf{构造达到上界的游向及时间比例}\par\nobreak
现在还须构造达到上界的方案：取游泳方向 $\bm e=\bm n$，即与岸成60°。由纵分量要求 $2\sqrt2\sin15^\circ=\sqrt3(1-\lambda)$。先用差角公式计算
\[
 \sin15^\circ=\sin(45^\circ-30^\circ)
 =\frac{\sqrt6-\sqrt2}{4},\qquad
 2\sqrt2\sin15^\circ=\sqrt3-1.
\]
因此 $\sqrt3(1-\lambda)=\sqrt3-1$，移项得 $\sqrt3\lambda=1$，即 $\lambda=1/\sqrt3$。
再检验横分量：$4\lambda+(1-\lambda)=1+\sqrt3=2\sqrt2\cos15^\circ$，也恰好满足，且 $0<\lambda<1$。所以给出的跑、游阶段都具有正时间，确实能到同一相遇点。最终
\[\boxed{U_{\max}=2\sqrt2\ \mathrm{m/s}},\qquad
\boxed{t_{\text{跑}}:t_{\text{游}}=1:(\sqrt3-1)}.
\]
这一比例与总时间无关；题目没有限定相遇距离时，任取正总时间即可按比例实现。

投影不等式证明所有方案都不能更快，分量检验则证明这个上界能够达到。\end{drillsolution}

\begin{drillsolution}{k6p18}\textbf{由对称规则求边长缩短率和相遇时间}\par\nobreak
由图\ref{k6p18q}的循环追逐规则及等边初始构形，三只蜗牛的运动在旋转120°并重新编号后保持不变，所以相遇前始终构成等边三角形，中心 $O$固定。记边长为 $a(t)$，爬速为 $u$，沿边 $A\to B$取单位矢量 $\bm e$。边长变化率是两端相对速度的投影：
\[\dot a=(\bm v_B-\bm v_A)\cdot\bm e.
\]
$A$正朝 $B$运动，投影为 $u$；$B$朝 $C$运动，与 $A\to B$成120°，投影为 $-u/2$。

因此 $\dot a=-3u/2$。用初值 $a(0)=a_0$ 积分：
\[
a(t)-a_0=-\frac{3u}{2}\int_0^t\dd s
=-\frac{3u}{2}[s]_0^t=-\frac{3ut}{2},
\]
即 $a=a_0-3ut/2$。令边长为零，得到
\[\boxed{T=\frac{2a_0}{3u}=8\text{分钟}},\qquad
\boxed{s=uT=40\text{厘米}}.
\]
\fig{k6p18s}{三只蜗牛的对称追逐：解答过程}{\begin{tikzpicture}\foreach \a in {90,210,330}\draw[blue,world]plot[domain=0:420,samples=120]({1.8*exp(-sqrt(3)*rad(\x))*cos(\a+\x)},{1.8*exp(-sqrt(3)*rad(\x))*sin(\a+\x)});\fill(0,0)circle(2pt)node[below]{$O$};\draw[flow,red](2.5,1)--(2.5,-.5)node[right]{$\sqrt3u/2$};\draw[flow,blue](2.5,1)--(3.4,1)node[above]{$u/2$};\node at(2.8,-1.1){径向与切向};\end{tikzpicture}}



\par\medskip\Needspace{4\baselineskip}\noindent\textbf{在径向与切向分解速度}\par\nobreak
为了求轨迹，按图\ref{k6p18s}将每只速度分成向心和切向两部分。到中心距离 $\rho=a/\sqrt3$，沿追逐方向取方位角增加，分量关系给
\[\dot\rho=-\frac{\sqrt3}2u,\qquad \rho\dot\varphi=\frac u2.
\]


\par\medskip\Needspace{4\baselineskip}\noindent\textbf{积分轨迹并检查有限时间与无限转角}\par\nobreak
相除消去时间得到 $\dd\rho/\dd\varphi=-\sqrt3\rho$。相遇前 $\rho>0$，分离变量，并从初始状态作定积分：
\[
\int_{\rho_0}^{\rho}\frac{\dd s}{s}
=-\sqrt3\int_{\varphi_0}^{\varphi}\dd\psi,\qquad
[\ln s]_{\rho_0}^{\rho}=-\sqrt3[\psi]_{\varphi_0}^{\varphi}.
\]
代入上下限，再取指数，得到
\[
\ln\frac\rho{\rho_0}=-\sqrt3(\varphi-\varphi_0),\qquad
\boxed{\rho=\rho_0e^{-\sqrt3(\varphi-\varphi_0)}}.
\]
故轨迹为对数螺线。由 $\rho=\rho_0(1-t/T)$ 和 $\dot\varphi=u/(2\rho)$，
令 $w=1-s/T$，$\dd s=-T\,\dd w$，从时间方程也可直接积分：
\begin{align*}
\varphi-\varphi_0
&=\frac{u}{2\rho_0}\int_0^t\frac{\dd s}{1-s/T}
=-\frac{uT}{2\rho_0}\int_1^{1-t/T}\frac{\dd w}{w}\\
&=-\frac1{\sqrt3}[\ln w]_1^{1-t/T}
=-\frac1{\sqrt3}\ln(1-t/T).
\end{align*}
这说明 $t\to T^-$时转角无界，但 $\rho\to0$、剩余路程也趋零。

路程40cm大于起点到中心的直线距离 $60/\sqrt3$cm，小于初始边长60cm，符合弯曲追逐而非直接奔向中心的图形。\end{drillsolution}
